\textbf{A shared acquired state develops different spatial futures under continued and stopped writing.}
This figure preserves the archived core7501 true fields, full native chart and native update schematic. The continued-writing trajectory is reused from job \texttt{9bfb1b8c-a7b6-4c0f-a4d5-ae5073e26203}; the no-new-writing continuation is from job \texttt{72ceda07-54a9-47ed-ae54-03d954d78ee9}. Both futures start from the identical acquired step-240 body, internal fields and remaining world. The latter intervention sets the writing rate to zero at the fork; it is neither life-zero memory removal nor freezing of memory. No scientific trajectory was recomputed for this figure.

\textbf{a: spatial divergence.} The shared state and the two futures at steps 400, 560 and 720 use the original periodic $32\times32$ coordinates, with no recentering, rotation, rescaling between branches or new field rendering. Density uses the preserved square-root colour scale from 0 to 0.75, with the original saturation. The scale bar spans 8 cells. Gold outlines mark food sites and coral outlines mark hazard sites. F and H identify the Food and Hazard windows enlarged in b. In this figure, F is a food-site identifier, not the fuel stock F used in Fig.~2. Local accumulation and displacement of a pattern do not by themselves establish coherent whole-body migration.

\textbf{b: local carried state.} Cue-channel $m_0$ at the food window and $m_1$ at the hazard window are shown at step 400 on the preserved $-1$ to $+1$ scale. White indicates near-zero $m$; grey masks $\rho<0.001$. These local states are measured along the actual diverging trajectories. They are not independently isolated mediators, and their visible association with later motion does not establish a mediated fraction.

\textbf{c: specific-energy field.} The step-560 field is $E/\rho$, in EU/M, on the preserved scale 0--0.16, saturated above 0.16. White is near 0.08; grey again masks $\rho<0.001$. The values 56.23\% and 16.85\% are fractions of all tissue mass below 0.08 EU/M, including masked cells. The specific-energy field is an observed state, not an isolated causal mediator or an independent fitness measure.

\textbf{d: complete saved-time comparison.} The native editable chart retains both complete 61-point series from 240 to 720. At 560 the displayed masses are 20.72 and 16.13 M; at 720 they are 5.90 and 8.72 M for continued and stopped writing, respectively. Remaining energy has the opposite endpoint ordering, 4.755 versus 1.170 EU. Over 240--720, continued-minus-stopped injury loss is $-6.062$ M and starvation loss is $+8.889$ M. These are recorded stock/flow differences, not isolated causal pathways or mediation fractions. Saved times are not independent repeats.

\textbf{e: update ordering.} Separate symbolic rows illustrate a local food-contact episode, with circles representing carried memory; arrows are not measured fluxes or drawn to scale. The current transport map carries tissue and internal loads before cue-driven eligibility gates memory writing. Memory then contributes to subsequent desired motion, which is limited by available donor energy before actual transport. Turning off the new writing term retains eligibility dynamics, passive memory decay, memory transport and autonomous polarity dynamics. The diagram is an operator description, not a tested mediation chain.

The frozen transport model has no tissue-birth term: conserved movement redistributes existing tissue, after which injury and starvation remove tissue and carried energy. Uptake can replenish energy, not create tissue in this model. This spatial writing comparison is therefore distinct from the paper's paid generative and tissue-replacement models. The two futures are one common-state comparison in a development world, not independent confirmation or a general survival/learning advantage.
